\section{Foundation Knowledge}
\label{sec:foundation}

\subsection{Recommendation Systems and Sequential Modelling}
\label{sec:seq_rec}

A recommendation system aims to predict items that a user is likely to interact with, given their historical behaviour.
Three broad paradigms have been studied in the literature.
\textit{Collaborative filtering} constructs a user–item interaction matrix and exploits co-occurrence patterns, typically via matrix factorisation, to infer latent user and item factors.
\textit{Content-based filtering} represents items by their attributes (title, genre, description) and retrieves items similar to those a user has previously consumed.
\textit{Hybrid} approaches combine both signals to mitigate the individual weaknesses of each: collaborative methods suffer from the cold-start problem when a new item has no interactions, while content-based methods ignore the community-level purchasing signals that often drive discovery.

Sequential recommendation refines the problem further by treating a user's interactions as an ordered sequence and predicting the \textit{next} item.
The key observation is that a user's near-term intent is more accurately reflected by their most recent interactions than by their full historical profile.
Early sequential methods applied recurrent neural networks (RNNs) to model temporal dynamics~\cite{hidasi2016gru4rec}, but RNNs encode history into a fixed-size state vector, creating an information bottleneck for long sequences.

Self-attention, introduced in the Transformer architecture~\cite{vaswani2017attention}, eliminates this bottleneck by computing a weighted sum over all positions in the sequence for each query position.
The attention weight between positions $i$ and $j$ is:
\begin{equation}
  \alpha_{ij} = \frac{\exp((\mathbf{q}_i \mathbf{k}_j^\top) / \sqrt{d})}{\sum_{l} \exp((\mathbf{q}_i \mathbf{k}_l^\top) / \sqrt{d})}
\end{equation}
where $\mathbf{q}_i$, $\mathbf{k}_j$ are query and key projections of the $i$-th and $j$-th sequence positions, and $d$ is the head dimension.
SASRec~\cite{kang2018self} adapts this to recommendation with a \textit{causal} (left-to-right) mask so that only past interactions influence each prediction, matching the sequential nature of the task.
BERT4Rec~\cite{sun2019bert4rec} instead uses bidirectional attention with masked-item training, achieving higher accuracy but at a greater inference cost.

DIF-SASRec~\cite{xie2022difsasrec} extends SASRec by incorporating \textit{side information} — item attributes such as textual content and categorical metadata — through decoupled attention streams.
Rather than fusing attribute embeddings into item ID embeddings before attention (which conflates the two signals), DIF-SASRec maintains separate attention modules for content and category features and combines them at the output layer.
This decoupling preserves the ability to attend to content-relevant history and category-relevant history independently, improving performance on items with rich metadata.

\subsection{Graph-Based Behavioural Embedding: Cleora}
\label{sec:cleora}

Graph-based embedding methods represent items as nodes in a graph and learn dense vector representations that encode the graph's structural proximity.
In a co-purchase graph, items are connected by an edge if they are frequently bought together; the edge weight reflects the strength of that co-occurrence.
Dense items in tightly connected communities — such as books in the same series or by the same author — receive similar embeddings, enabling efficient candidate retrieval by nearest-neighbour search.

Cleora~\cite{rychalska2021cleora} is an unsupervised graph embedding scheme based on iterative feature propagation.
Starting from random node features, each iteration replaces every node's embedding with the weighted average of its neighbours' embeddings, followed by row-wise $\ell_2$ normalisation.
Formally, for node $v$ at iteration $t$:
\begin{equation}
  \mathbf{e}_v^{(t)} = \frac{\sum_{u \in \mathcal{N}(v)} w_{uv}\, \mathbf{e}_u^{(t-1)}}{\left\| \sum_{u \in \mathcal{N}(v)} w_{uv}\, \mathbf{e}_u^{(t-1)} \right\|_2}
\end{equation}
where $\mathcal{N}(v)$ is the neighbour set of $v$ and $w_{uv}$ is the edge weight.
Unlike DeepWalk or Node2Vec, Cleora requires no random-walk sampling and no contrastive objective, making it substantially faster and scalable to hundreds of thousands of nodes on a single CPU.
After a fixed number of iterations (typically two to three), the resulting embeddings are compact representations of co-purchase community membership.

\subsection{Dense and Sparse Retrieval with Hybrid Fusion}
\label{sec:retrieval}

\subsubsection{Bi-Encoder Dense Retrieval}

Dense retrieval represents both queries and documents as fixed-length vectors and retrieves candidates by computing the inner product (or cosine similarity) between query and document embeddings.
A \textit{bi-encoder} architecture uses two independent encoders — one for queries and one for items — allowing item vectors to be pre-computed and indexed offline.
At query time, only the query encoding and a vector lookup are needed, making dense retrieval highly efficient.

BGE-M3~\cite{chen2024bge} (\texttt{BAAI/bge-m3}) is the encoder used in this system.
It supports 100+ languages and produces 1{,}024-dimensional dense embeddings.
Its multilingual pre-training via self-knowledge distillation enables strong retrieval quality on both English and Vietnamese queries without language-specific fine-tuning.

\subsubsection{Approximate Nearest Neighbour Search}

Exact nearest-neighbour search over millions of vectors scales as $O(n \cdot d)$ per query, where $d$ is the embedding dimension.
For $n = 1{,}700{,}000$ vectors of dimension 1{,}024 this is prohibitively slow at interactive latencies.
Hierarchical Navigable Small World (HNSW)~\cite{malkov2018hnsw} graphs provide approximate nearest-neighbour search in $O(\log n)$ expected time.
HNSW builds a layered proximity graph: lower layers are dense with short-range edges and upper layers are sparse with long-range edges.
At query time, a greedy search descends from the sparser upper layers to progressively refine the candidate set, allowing large portions of the index to be skipped.
FAISS~\cite{johnson2021faiss} provides a GPU-accelerated implementation of both HNSW and flat (exact) indices used in this system.

\subsubsection{Sparse Retrieval and Reciprocal Rank Fusion}

BM25 is a probabilistic sparse retrieval model that scores a document $d$ for query $q$ as:
\begin{equation}
  \text{BM25}(q, d) = \sum_{t \in q} \text{IDF}(t) \cdot \frac{f(t, d) \cdot (k_1 + 1)}{f(t, d) + k_1 \left(1 - b + b \cdot \frac{|d|}{\text{avgdl}}\right)}
\end{equation}
where $f(t, d)$ is the term frequency of token $t$ in document $d$, $|d|$ is document length, $\text{avgdl}$ is the corpus average document length, and $k_1 = 1.2$, $b = 0.75$ are standard hyperparameters.
Tantivy, a Rust-native full-text search library, provides BM25 indexing and retrieval over the item catalogue.

Dense and sparse rankings capture complementary signals — dense retrieval finds semantically similar items even when the query shares no tokens with the document, while BM25 excels at exact-match term recall.
Reciprocal Rank Fusion (RRF)~\cite{cormack2009rrf} combines two ranked lists without score normalisation:
\begin{equation}
  \text{RRF}(d) = \sum_{r \in \mathcal{R}} \frac{1}{k + r_i(d)}
\end{equation}
where $\mathcal{R}$ is the set of input rankings, $r_i(d)$ is the rank of document $d$ in ranking $i$, and $k = 60$ is a smoothing constant.
RRF is robust to score scale differences between retrieval systems and has been shown to outperform learned rank aggregation methods in low-resource settings~\cite{cormack2009rrf}.

\subsection{Multimodal and Multilingual Extensions}
\label{sec:multimodal}

\paragraph{Visual Search (CLIP) —}
CLIP~\cite{radford2021clip} is a vision-language model pre-trained via contrastive learning on 400 million image–text pairs.
It maps both images and text descriptions into a shared embedding space, enabling zero-shot image search: a query image can be encoded and matched against text-based item embeddings without domain-specific fine-tuning.
In the system, a CLIP HNSW index over 3 million item cover images enables cover-image-based book search.

\paragraph{Multilingual Query Handling (NLLB) —}
When a query arrives in Vietnamese, the NLLB-200 model~\cite{nllb2022} translates it to English before the BGE-M3 encoding step.
NLLB-200 is a 54-billion-parameter mixture-of-experts model supporting 202 languages, trained specifically to improve translation quality for low-resource languages.
Query-time translation enables the system to leverage the full semantic richness of the BGE-M3 English-dominant training distribution while still accepting Vietnamese input natively.

